\documentclass[border=4pt]{standalone}
\usepackage{amsmath,amssymb,mathtools,xcolor,varwidth}
\begin{document}
\begin{varwidth}{28em}
\large\bfseries
Below sits the second figure $PprT$ with its own rank of three parallelograms on the equal bases $PQ$, $QR$, $RT$ of the baseline $PT$. Notice that the curve $prT$ leans the other way --- it descends from $p$ toward $T$ --- so this figure is shaped quite differently from the first. Yet it carries the \emph{same count} of strips, which we name $q_1$, $q_2$, $q_3$ in order from $P$ toward $T$. The whole point of Lemma IV is that the figures may differ wildly in form, provided their corresponding strips agree in ultimate ratio.
\end{varwidth}
\end{document}
